
\section{Conclusion}
\label{sec:conclusion}


Qwen2.5 represents a significant advancement in large language models (LLMs), with enhanced pre-training on 18 trillion tokens and sophisticated post-training techniques, including supervised fine-tuning and multi-stage reinforcement learning. These improvements boost human preference alignment, long text generation, and structural data analysis, making Qwen2.5 highly effective for instruction-following tasks.
Available in various configurations, Qwen2.5 offers both open-weight from 0.5B to 72B parameters and proprietary models including cost-effective MoE variants like Qwen2.5-Turbo and Qwen2.5-Plus. Empirical evaluations show that Qwen2.5-72B-Instruct matches the performance of the state-of-the-art Llama-3-405B-Instruct, despite being six times smaller. Qwen2.5 also serves as a foundation for specialized models, demonstrating its versatility for domain-specific applications.
We believe that Qwen2.5's robust performance, flexible architecture, and broad availability make it a valuable resource for both academic research and industrial applications, positioning it as a key player of future innovations.


In the future, we will focus on advancing robust foundational models. 
First, we will iteratively refine both base and instruction-tuned large language models (LLMs) by incorporating broader, more diverse, higher-quality data.
Second, we will also continue to develop multimodal models. Our goal is to integrate various modalities into a unified framework. This will facilitate seamless, end-to-end information processing across textual, visual, and auditory domains.
Third, we are committed to enhancing the reasoning capabilities of our models. This will be achieved through strategic scaling of inference compute resources. 
These efforts aim to push the boundaries of current technological limitations and contribute to the broader field of artificial intelligence.
